\documentclass[11pt,letterpaper]{article}

% --------------------------------------------------
% Packages
% --------------------------------------------------
\usepackage[T1]{fontenc}
\usepackage[utf8]{inputenc}
\usepackage{lmodern}
\usepackage{microtype}
\usepackage[margin=1in]{geometry}
\usepackage{amsmath}
\usepackage{amssymb}
\usepackage{booktabs}
\usepackage{tabularx}
\usepackage{array}
\usepackage{enumitem}
\usepackage[hidelinks]{hyperref}

% --------------------------------------------------
% Formatting
% --------------------------------------------------
\setlength{\parindent}{0pt}
\setlength{\parskip}{0.7em}
\renewcommand{\arraystretch}{1.25}

\newcolumntype{Y}{>{\raggedright\arraybackslash}X}

\setlist[itemize]{
    leftmargin=*,
    itemsep=0.2em,
    topsep=0.3em
}
\setlist[enumerate]{
    leftmargin=*,
    itemsep=0.4em,
    topsep=0.3em
}

% --------------------------------------------------
% Document
% --------------------------------------------------
\title{
    \textbf{Proposal: 2.4 GHz Low-Power\\
    Gilbert-Cell Downconversion Mixer}
}
\author{
    Sangeon Jeon, Yunho Cho \\
    Chipalooza Challenge \#4\\
    SkyWater SKY130 CMOS
}
\date{November 2026}

\begin{document}

\maketitle

% ==================================================
\section{IP Block Type and Overview}
% ==================================================

This project proposes the design and implementation of a
\textbf{2.4 GHz low-power active downconversion mixer}
using the open-source SkyWater SKY130 CMOS process.

The proposed IP will use a double-balanced Gilbert-cell
architecture to translate RF signals in the 2.4 GHz ISM
band to a low intermediate frequency (IF). It is intended
as a reusable analog building block for wireless receivers,
including Wi-Fi, Bluetooth, IoT, and other RF communication
systems.

The primary objective is to provide a fully open-source,
silicon-validated mixer IP with documented conversion gain,
noise figure, linearity, port isolation, power consumption,
and input impedance characteristics.

The baseline design will use fixed biasing. If development
time permits, digitally programmable biasing will be added
as an optional enhancement to investigate
power--performance trade-offs.


% ==================================================
\section{Functional Description}
% ==================================================

The proposed mixer will consist of the following
subcircuits:

\begin{enumerate}
    \item \textbf{RF Transconductance Stage:}
    Converts the incoming RF voltage into a differential
    current for frequency conversion.

    \item \textbf{Gilbert-Cell Switching Core:}
    Performs frequency translation using a differential
    LO switching pair.

    \item \textbf{LO Input Interface:}
    Accepts an externally generated LO signal and
    provides the complementary signals required by the
    switching core.

    \item \textbf{IF Load and Output Interface:}
    Converts the downconverted current into an IF voltage
    suitable for measurement or connection to a
    subsequent analog stage.

    \item \textbf{Bias Circuitry:}
    Establishes the operating currents and bias points
    of the mixer.
\end{enumerate}

The mixer will operate according to

\begin{equation}
    f_{\mathrm{IF}}
    =
    \left|f_{\mathrm{RF}}-f_{\mathrm{LO}}\right|.
\end{equation}

For example, an RF input at \(2.450\,\mathrm{GHz}\) and an
LO input at \(2.449\,\mathrm{GHz}\) will produce a desired
IF component at \(1\,\mathrm{MHz}\).

% (?) 너무 낮지 않냐? 보통 ADC sampling rate

The design will target low power consumption while
maintaining useful conversion gain, noise performance,
and linearity.

Input impedance matching will be considered for
compatibility with standard \(50\,\Omega\) RF measurement
equipment. The matching network will be designed and
characterized together with the input circuitry and
relevant pad parasitics.


% ==================================================
\section{I/O and Interface Requirements}
% ==================================================

The baseline design requests three dedicated analog
signal pads.

\begin{table}[htbp]
\centering
\caption{Proposed I/O signals and interfaces}
\label{tab:io}
\begin{tabularx}{\textwidth}{@{}llY@{}}
\toprule
\textbf{Signal} &
\textbf{Direction} &
\textbf{Description} \\
\midrule

\texttt{RF\_IN} &
Input &
Single-ended RF input, nominally
\(2.4\)--\(2.5\,\mathrm{GHz}\) \\

\texttt{LO\_IN} &
Input &
Single-ended external LO input \\

\texttt{IF\_OUT} &
Output &
Buffered, single-ended low-frequency IF output \\

\texttt{VDD / GND} &
Power &
Supplied through the Chipalooza harness \\

\texttt{ENABLE} &
Digital input &
Mixer enable/power-down control, if available \\

\bottomrule
\end{tabularx}
\end{table}

The mixer will internally generate the complementary
signals necessary for Gilbert-cell operation.

\textbf{Pad requirements:}
\texttt{RF\_IN} and \texttt{LO\_IN} require dedicated
connections suitable for GHz-frequency signaling.
\texttt{IF\_OUT} also requests a dedicated pad to simplify
independent characterization. The final pad arrangement
and compatibility with the SKY130 pad library will be
confirmed during proposal and schematic review.

The baseline circuit will target a nominal
\(1.8\,\mathrm{V}\) core supply, subject to confirmation
of harness supply availability, pad voltage limits,
and required interface protection.

External RF and LO signal generators will be used for
characterization; an on-chip PLL or LO frequency
synthesizer is outside the scope of this proposal.


% ==================================================
\section{Preliminary Target Specifications}
% ==================================================

The following values are initial design targets rather
than demonstrated performance. Final feasibility will be
assessed through SKY130 transistor-level simulations and
review of the RF pad interface.

\begin{table}[htbp]
\centering
\caption{Preliminary mixer specifications}
\label{tab:specifications}
\begin{tabularx}{\textwidth}{@{}p{0.46\textwidth}Y@{}}
\toprule
\textbf{Parameter} & \textbf{Preliminary target} \\
\midrule

CMOS process &
SkyWater SKY130 \\

RF operating band &
\(2.4\)--\(2.5\,\mathrm{GHz}\) \\

Nominal RF frequency &
\(2.45\,\mathrm{GHz}\) \\

IF test frequency &
\(0.5\)--\(5\,\mathrm{MHz}\) \\

Core supply voltage &
\(1.8\,\mathrm{V}\) nominal \\

Supply characterization &
\(1.62\)--\(1.98\,\mathrm{V}\),
subject to PDK limits \\

Conversion gain &
\(\geq 5\,\mathrm{dB}\) nominal;
\(\geq 0\,\mathrm{dB}\) minimum target \\

Single-sideband noise figure &
\(\leq 15\,\mathrm{dB}\) nominal target \\

Input return loss (\(S_{11}\)) &
\(\leq -10\,\mathrm{dB}\)
at nominal RF frequency \\

Input third-order intercept (\(\mathrm{IIP3}\)) &
\(\geq -15\,\mathrm{dBm}\) nominal target \\

LO-to-RF isolation &
\(\geq 20\,\mathrm{dB}\) \\

DC power consumption &
\(\leq 20\,\mathrm{mW}\) nominal \\

Load for baseline IF testing &
High-impedance measurement input \\

\bottomrule
\end{tabularx}
\end{table}

Specifications will be evaluated under clearly defined
RF input power, LO drive, load impedance, supply voltage,
and temperature conditions. Transistor operating voltages
and RF input levels will remain within the applicable
PDK and pad absolute limits.

The design will be characterized across process, voltage,
and temperature corners, including post-layout extracted
parasitics. The acceptance limits and simulation methods
will be finalized during schematic review with the
challenge organizers.


% ==================================================
\section{Verification and Measurement Plan}
% ==================================================

The mixer will undergo pre-layout simulation, physical
implementation, post-layout verification, and silicon
characterization.

\subsection{Pre-Layout Verification}

\begin{itemize}
    \item Verify DC operating points and device
    operating regions.

    \item Evaluate conversion gain across RF frequency,
    LO drive, and bias conditions.

    \item Characterize RF input impedance and \(S_{11}\).

    \item Evaluate linearity using two-tone RF tests
    and estimate IIP3.

    \item Evaluate LO-to-RF feedthrough and port
    isolation.

    \item Analyze noise performance and power
    consumption.

    \item Perform process, voltage, temperature,
    and mismatch sensitivity analyses.
\end{itemize}


\subsection{Physical Design and Post-Layout Verification}

The mixer will be laid out using the SKY130 PDK, with
attention to differential symmetry, device matching,
parasitic coupling, and RF signal routing.

Verification will include:

\begin{itemize}
    \item Design Rule Check (DRC)

    \item Layout Versus Schematic (LVS)

    \item Parasitic extraction (PEX)

    \item Post-layout conversion gain, linearity,
    isolation, and power simulations

    \item PVT verification using reproducible
    open-source workflows
\end{itemize}


\subsection{Silicon Characterization}

Following fabrication, RF and LO signals will be
applied through the dedicated test pads.

The IF output will be measured using appropriate test
equipment to characterize conversion gain, operating
frequency range, and linearity.

A two-tone RF measurement will be used to determine
third-order intermodulation performance. Input matching
and LO leakage will be characterized using RF measurement
equipment where available.

Noise figure will be characterized using a calibrated
noise measurement method, with test-fixture and interface
losses accounted for.

Measurements will be compared with pre-layout and
post-layout simulation results, and any discrepancies
will be documented.


% ==================================================
\section{Optional Enhancement: Programmable Bias}
% ==================================================

If the baseline mixer meets its target specifications
and the design schedule permits, programmable bias
circuitry will be incorporated.

This enhancement will allow selection among several
operating bias points to investigate trade-offs between:

\begin{itemize}
    \item Conversion gain
    \item Noise figure
    \item Third-order linearity
    \item DC power consumption
\end{itemize}

The implementation may use a digitally controlled
current mirror or a switchable bias network.

\textbf{Programmable bias is not required for successful
operation or verification of the baseline mixer.}
The fixed-bias mixer will remain a standalone,
reusable IP block.


% ==================================================
\section{Open-Source Deliverables}
% ==================================================

The project will be released under the Apache-2.0
license, subject to compatibility with all incorporated
source materials.

The public repository will include:

\begin{itemize}
    \item Transistor-level schematics and simulation
    testbenches

    \item Circuit documentation and design
    specifications

    \item Standalone layout and Chipalooza harness
    integration views

    \item DRC/LVS reports and parasitic extraction
    results

    \item Reproducible scripts for open-source
    EDA verification

    \item Post-layout simulation summaries

    \item Silicon measurement results following
    fabrication
\end{itemize}

The design will be developed and verified using the
SKY130 open-source PDK, with reproducible verification
through tools such as Xschem, ngspice, Magic, and/or
KLayout as appropriate.


% ==================================================
\section{Project Scope and Schedule}
% ==================================================

The project will prioritize a complete baseline
Gilbert-cell mixer over optional functionality.

\begin{table}[htbp]
\centering
\caption{Proposed project milestones}
\label{tab:milestones}
\begin{tabularx}{\textwidth}{@{}p{0.43\textwidth}Y@{}}
\toprule
\textbf{Milestone} & \textbf{Deliverable} \\
\midrule

Proposal review --- Nov.\ 2026 &
Architecture, target specifications, and I/O plan \\

Schematic review --- Dec.\ 2026 &
Complete circuit and pre-layout verification \\

Layout review --- Feb.\ 2027 &
Standalone layout, DRC/LVS, PEX and post-layout
verification \\

Final review --- Feb.\ 2027 &
Harness integration and sign-off \\

Tapeout --- May 2027 &
Final open-source design submission \\

\bottomrule
\end{tabularx}
\end{table}

The proposed scope is deliberately limited to a
single RF downconversion mixer. A complete receiver,
on-chip PLL, and automatic I/Q calibration are not
included.


% ==================================================
\section{Expected Contribution}
% ==================================================

The project aims to provide an accessible, reusable,
and silicon-characterized 2.4 GHz analog mixer IP
in SKY130 CMOS.

Although the Gilbert-cell topology is well established,
an openly documented implementation with reproducible
verification and measured RF performance would provide
a practical building block for educational, research,
and open-source wireless SoC development.

The optional programmable-bias extension may further
demonstrate how conversion gain, linearity, noise, and
power consumption vary with operating conditions,
helping future designers select appropriate
configurations for their applications.

\end{document}


% Timeline
% - Wednesday 11/18 : Draft (Yunho)

% Sections
% - Team
% - Overview/Description
% - Specification ← 어케 함 ??? LLM! refer to final project _2026_Fall.pdf
        % S11
        % Noise Factor
        % IIP3
        % Tempurature range
        % bias current
        % IRR
% - Pinout
% (- Architecture) 는 너무 이를듯
% - External resources
% - Challenges
% - Test Plan

% Decisions to make:
% - Programmability for bias current?   

\section{Introduction}


\end{document}

